\selectlanguage{spanish}
\usepackage[utf8]{inputenc}
\usepackage[T1]{fontenc}
\usepackage[spanish,es-tabla,es-nodecimaldot]{babel}
\usepackage{lmodern}
\addto\extrasspanish{\spanishucroman}
\spanishucroman
\usepackage{microtype}
\usepackage{amsmath,amssymb}
\usepackage{graphicx}
\usepackage{tikz}
\usetikzlibrary{arrows.meta,babel}
\usepackage{float}
\usepackage{booktabs}
\usepackage{tabularx}
\usepackage{longtable}
\usepackage{etoolbox}
\usepackage{array}
\usepackage{makecell}
\usepackage{multirow}
\usepackage{enumitem}
\usepackage[table]{xcolor}
\usepackage{caption}
\usepackage{subcaption}
\usepackage{siunitx}
\usepackage{url}
\usepackage{xurl}
\usepackage{hyperref}
\usepackage{fancyhdr}
\usepackage{setspace}
\usepackage{csquotes}
\usepackage[backend=biber,style=ieee,sorting=none,defernumbers=true,maxbibnames=99]{biblatex}
\addbibresource{bibliografia/referencias.bib}
\addbibresource{aspersion_manual/referencias.bib}
\AtBeginBibliography{\raggedright}
% Formato IEEE: sin comillas angulares y con la formula de acceso en linea del estilo IEEE.
\DefineBibliographyStrings{spanish}{urlseen={Accessed:},url={[Online]. Available}}
\AtBeginDocument{\setquotestyle{english}}

\usepackage{newunicodechar}
\newunicodechar{≤}{\ensuremath{\leq}}
\newunicodechar{≥}{\ensuremath{\geq}}
\newunicodechar{≈}{\ensuremath{\approx}}
\newunicodechar{−}{\ensuremath{-}}
\newunicodechar{×}{\ensuremath{\times}}
\newunicodechar{±}{\ensuremath{\pm}}
\newunicodechar{Ø}{\O{}}

\graphicspath{{images/}{images/institucional/}{images/mesa/}{images/sap/}{images/aspersion/}{images/gestion/}{images/p3/}{images/planos/}}

\definecolor{utecblue}{RGB}{0,78,121}
\definecolor{statusgreen}{RGB}{25,115,70}
\definecolor{statusamber}{RGB}{155,95,0}
\definecolor{statusred}{RGB}{160,35,35}
\definecolor{lightgray}{RGB}{242,244,246}

\hypersetup{
  colorlinks=true,
  linkcolor=black,
  citecolor=black,
  urlcolor=black,
  filecolor=black,
  pdftitle={Diseño de una máquina para la limpieza de rodamientos HPGR},
  pdfauthor={Ismael Enrique Rodriguez Rojas y Jean Piere Fredes Arapa},
  pdfsubject={Informe Capstone de Ingeniería Mecánica}
}

\sisetup{output-decimal-marker={,},group-separator={\,},group-minimum-digits=4}
\onehalfspacing
\setlength{\emergencystretch}{3em}
\setlength{\LTpre}{0.8em}
\setlength{\LTpost}{0.8em}
\setlist[itemize]{topsep=3pt,itemsep=2pt,parsep=0pt}
\setlist[enumerate]{topsep=3pt,itemsep=3pt,parsep=0pt}
\captionsetup{font=small,labelfont=bf,justification=centering}
\cftsetindents{chapter}{0em}{9.5em}
\cftsetindents{section}{1.5em}{3em}
\cftsetindents{subsection}{3em}{4em}
\cftsetindents{figure}{0em}{6em}
\cftsetindents{table}{0em}{6em}
\renewcommand{\cftchapfont}{\bfseries}
\renewcommand{\cftchappresnum}{CAPÍTULO\space}
\renewcommand{\cftfigfont}{\normalfont}
\renewcommand{\cftfigpresnum}{Figura\space}
\renewcommand{\cfttabfont}{\normalfont}
\renewcommand{\cfttabpresnum}{Tabla\space}
\setlength{\cftbeforechapskip}{0pt}
\setlength{\cftbeforesecskip}{0pt}
\setlength{\cftbeforesubsecskip}{0pt}
\setlength{\cftbeforefigskip}{0pt}
\setlength{\cftbeforetabskip}{0pt}

\makeatletter
\newcommand{\indicepreliminar}[2]{%
  \cleardoublepage
  \begingroup
  \doublespacing
  \setlength{\parskip}{0pt}
  \chapter*{#1}
  \let\addvspace\@gobble
  \@starttoc{#2}
  \cleardoublepage
  \endgroup
}
\renewcommand{\tableofcontents}{\indicepreliminar{TABLA DE CONTENIDO}{toc}}
\renewcommand{\listoftables}{\indicepreliminar{ÍNDICE DE TABLAS}{lot}}
\renewcommand{\listoffigures}{\indicepreliminar{ÍNDICE DE FIGURAS}{lof}}
\makeatother

\newcommand{\preliminarpendiente}[2]{%
  \cleardoublepage
  \thispagestyle{empty}
  \phantomsection
  \addcontentsline{toc}{chapter}{#1}
  \null\vfill
  \noindent\makebox[\textwidth][r]{\normalfont\normalsize\itshape #2:}\par
  \vfill\null
  \cleardoublepage
}

\pagestyle{fancy}
\fancyhf{}
\fancyhead{}
\fancyfoot[C]{\thepage}
\setlength{\headheight}{16pt}

\newcommand{\estadoverificado}{{\scriptsize\textcolor{black}{\textbf{VERIFICADO}}}}
\newcommand{\estadocalculado}{{\scriptsize\textcolor{black}{\textbf{CALCULADO}}}}
\newcommand{\estadoprovisional}{{\scriptsize\textcolor{black}{\textbf{PROVISIONAL}}}}
\newcommand{\estadopendiente}{{\scriptsize\textcolor{black}{\textbf{PENDIENTE}}}}
\newcommand{\estadoantecedente}{{\scriptsize\textcolor{black}{\textbf{ANTECEDENTE}}}}
\newcommand{\R}{\mathbb{R}}
\newcommand{\source}[1]{\par\smallskip{\footnotesize\itshape Procedencia: #1}\par}
\newcommand{\longstatus}[1]{\begingroup\footnotesize\path{#1}\endgroup}

\newcolumntype{P}[1]{>{\raggedright\arraybackslash\hspace{0pt}}p{#1}}
\newcolumntype{Y}{>{\raggedright\arraybackslash}X}
\newcolumntype{C}[1]{>{\centering\arraybackslash}p{#1}}
\newcolumntype{R}[1]{>{\raggedleft\arraybackslash}p{#1}}

\newenvironment{notaalcance}{%
  \begin{center}\begin{minipage}{0.94\textwidth}\small\color{black}\textbf{Alcance documental.}\ }%
  {\end{minipage}\end{center}}

\renewcommand{\chaptername}{CAPÍTULO}

\makeatletter
\patchcmd{\LT@output}{\vss}{\vfil}{}{}
\patchcmd{\LT@output}{\vss}{\vfil}{}{}
\makeatother
